\documentclass{article}

\usepackage[numbers]{natbib}
\usepackage[utf8]{inputenc}
\usepackage[T1]{fontenc}
\usepackage{hyperref}
\usepackage{url}
\usepackage{booktabs}
\usepackage{amsfonts}
\usepackage{caption}
\usepackage{subcaption}
\usepackage{amsmath}
\usepackage{amssymb}
\usepackage{graphicx}

\title{Learning When to Think: RL for Adaptive Reasoning Computation in LLMs}

\author{%
  Ivan Andhika \\
  School of Computing\\
  University of Utah\\
  \texttt{u1590903@utah.edu}
  \and
  Hao Ren \\
  School of Computing\\
  University of Utah\\
  \texttt{u1527543@utah.edu}
  \and
  Ammon Gleason \\
  School of Computing\\
  University of Utah\\
  \texttt{u1221580@utah.edu}
}

\begin{document}

\maketitle

\begin{abstract}
Large language models (LLMs) benefit from extended test-time computation, yet static reasoning budgets lead to computational waste on trivial queries and failure on complex ones.
While step-level verification improves reasoning quality, current reinforcement learning (RL) methods often suffer from ``mode collapse'' or ``overthinking'' due to uniform length penalties.
We propose a multi-turn adaptive policy with three specialized actions: \texttt{continue}, \texttt{refine}, and \texttt{terminate}.
Unlike passive verification, the \texttt{refine} action enables the model to utilize internal signals to correct intermediate errors.
We train a LoRA-adapted Qwen2.5-Math-7B-Instruct using Decoupled Group Relative Policy Optimization (DeGRPO) to balance mode-selection gradients against response accuracy.
To manage the overthinking dilemma, we implement an Adaptive Length Penalty (ALP) that scales inversely with problem difficulty.
We evaluate our approach on MATH-500 and AIME 2024 to provide sufficient headroom for policy discovery, comparing against fixed-compute and vanilla GRPO baselines.
\end{abstract}

\section{Introduction}

Increasing test-time computation through Chain-of-Thought (CoT) and self-consistency decoding has established a clear scaling law for reasoning.
However, recent evidence suggests that the efficiency of this scaling depends on instance-specific allocation: easy problems require minimal tokens, while difficult problems necessitate extended reflection or backtracking.
A core challenge in training adaptive models is ``mode collapse,'' where the model ignores control tokens in favor of standard reasoning, often triggered by gradient imbalances in standard RL algorithms.
Furthermore, static length penalties frequently lead to ``training collapse,'' where the model prematurely stops reasoning to avoid penalties.
We address these gaps by introducing: (1) a verify-and-refine MDP framework that treats error recovery as a high-reward skill; (2) an Adaptive Length Penalty (ALP) that redistributes the token budget based on the online solve rate; and (3) a DeGRPO training pipeline that decouples the mode-selection loss from the response content loss.

Related work on test-time scaling~\citep{wei2022chainofthought,wang2023selfconsistency,snell2024scaling}, process supervision~\citep{lightman2023lets,zhang2025lessons}, RL-style self-correction~\citep{kumar2024score}, and group-relative policy optimization in math models~\citep{shao2024deepseekmath} motivates our design choices.

\section{Method}

\subsection{MDP Formulation}

We model the reasoning process as a Markov Decision Process.
Given a query $q$, the state $s_t$ consists of the problem text and the trajectory generated so far.
The action space is defined as $\mathcal{A} = \{\texttt{continue}, \texttt{refine}, \texttt{terminate}\}$.

\begin{itemize}
  \item \textbf{Continue:} Extends the current reasoning chain.
  \item \textbf{Refine:} Triggers a self-correction block, specifically rewarded for ``progress'' (moving from an incorrect state to a correct one).
  \item \textbf{Terminate:} Concludes the sequence and outputs the final answer.
\end{itemize}

\subsection{Adaptive Length Penalty (ALP)}

To prevent overthinking on easy tasks without hindering complex reasoning, we employ an Adaptive Length Penalty.
The reward $R$ for a rollout $o$ is:
\begin{equation}
  R(o) = r_{\text{acc}} - \beta \cdot \max(0, \text{SR}(q)) \cdot \frac{n_{\text{tokens}}}{L_{\text{max}}}
\end{equation}
where $r_{\text{acc}} \in \{0, 1\}$ is the verifiable accuracy reward, $\text{SR}(q)$ is the online solve rate (empirical difficulty estimate) across rollouts in the group, and $\beta$ is the global penalty coefficient.

\subsection{Optimization via DeGRPO}

To stabilize training, we use Decoupled GRPO (DeGRPO), which separates the gradient updates for the initial control token and subsequent response tokens.
This prevents the high volume of response tokens from drowning out the signal for mode selection.
We apply LoRA to all linear layers (Attention and MLP modules), as MLP layers are critical for retrieving factual knowledge and complex logic.

\section{Experimental Design}

\subsection{Headroom and Benchmarks}

We move beyond the saturated GSM8K benchmark ($\sim$93\% baseline for our model) to MATH-500 and AIME 2024 ($\sim$16.6\% baseline).
This provides the necessary headroom to observe if RL discovers novel recovery strategies.

\subsection{Hypotheses}

\begin{enumerate}
  \item \textbf{H1:} DeGRPO with ALP achieves a superior accuracy-efficiency Pareto frontier than vanilla GRPO with static penalties.
  \item \textbf{H2:} The \texttt{refine} policy specifically improves performance on AIME by enabling intrinsic self-correction.
\end{enumerate}

\section{Limitations}

This scope omits Tool-Integrated Reasoning (TIR) and transfer to non-math domains like HumanEval.
Future work could integrate Qwen2.5-Math-PRM-7B for dense, step-level process rewards to accelerate convergence.

\bibliographystyle{plainnat}
\bibliography{sample}

\end{document}
